\section{Problem setting and normalization dictionary}
\label{sec:problem}

\subsection{Graph and PageRank objective}

Let \(G=(V,E)\) be a finite, connected, undirected, unweighted graph with no
isolated vertices.  Let \(A\) be its adjacency matrix,
\(D=\diag(d_1,\ldots,d_n)\) its degree matrix, and
\begin{equation}
    W:=D^{-1/2}AD^{-1/2}
    \label{eq:normalized-adjacency}
\end{equation}
its symmetric normalized adjacency matrix.  We fix a source vertex
\(s\in V\), write \(e_s\) for its coordinate vector, and take
\(\alpha\in(0,1/2)\) unless stated otherwise.

The normalization used by AESP is
\begin{equation}
    Q
    :=\frac{1+\alpha}{2}I-\frac{1-\alpha}{2}W,
    \qquad
    f(x)
    :=\frac12 x^\mathsf{T}Qx
       -\alpha x^\mathsf{T}D^{-1/2}e_s.
    \label{eq:aesp-objective}
\end{equation}
The spectrum satisfies
\begin{equation}
    \alpha I\preceq Q\preceq I,
    \label{eq:q-spectrum}
\end{equation}
so \(f\) is \(\alpha\)-strongly convex and \(1\)-smooth.  Its unique minimizer
is
\begin{equation}
    x^\star=\alpha Q^{-1}D^{-1/2}e_s,
    \qquad
    \pi=D^{1/2}x^\star.
    \label{eq:ppr-solution}
\end{equation}
The gradient is
\begin{equation}
    g(x):=\grad f(x)=Qx-\alpha D^{-1/2}e_s.
    \label{eq:original-gradient}
\end{equation}

The source result of \citet{huang2025accelerated} gives the coordinatewise
certificate
\begin{equation}
    \norm{D^{-1/2}g(x)}_\infty<\alpha\eps
    \quad\Longrightarrow\quad
    \norm{D^{-1}(D^{1/2}x-\pi)}_\infty<\eps.
    \label{eq:ppr-certificate}
\end{equation}
Thus a vertex is active at final tolerance \(\eps\) when
\begin{equation}
    \abs{g_u(x)}\geq \alpha\eps\sqrt{d_u}.
    \label{eq:final-active-unscaled}
\end{equation}
If the implementation stores the scaled gradient
\begin{equation}
    q(x):=D^{1/2}g(x),
    \label{eq:q-scaled-gradient}
\end{equation}
then the equivalent activation rule is
\begin{equation}
    \abs{q_u(x)}\geq\alpha\eps d_u.
    \label{eq:final-active-scaled}
\end{equation}
This distinction is purely notational but must be preserved in code and in the
work accounting.

\subsection{Dictionary to the locally evolving set normalization}

The NeurIPS 2024 locally evolving set paper uses the symmetric system
\begin{equation}
    \widehat Qx=\widehat b,
    \qquad
    \widehat Q:=I-\chi W,
    \qquad
    \chi:=\frac{1-\alpha}{1+\alpha},
    \qquad
    \widehat b:=\frac{2\alpha}{1+\alpha}D^{-1/2}e_s.
    \label{eq:work-one-system}
\end{equation}
The two normalizations are related by
\begin{equation}
    Q=\frac{1+\alpha}{2}\widehat Q,
    \qquad
    \alpha D^{-1/2}e_s
      =\frac{1+\alpha}{2}\widehat b.
    \label{eq:normalization-dictionary}
\end{equation}
For the Work-I residual
\begin{equation}
    \widehat r(x):=\widehat b-\widehat Qx,
    \label{eq:work-one-residual}
\end{equation}
we therefore have the exact identity
\begin{equation}
    g(x)=-\frac{1+\alpha}{2}\widehat r(x).
    \label{eq:gradient-residual-dictionary}
\end{equation}
Consequently, the AESP certificate in \cref{eq:ppr-certificate} is equivalent
to
\begin{equation}
    \norm{D^{-1/2}\widehat r(x)}_\infty
    <\frac{2\alpha}{1+\alpha}\eps,
    \label{eq:work-one-certificate}
\end{equation}
which is the stopping scale used by the locally evolving set solvers.

\begin{remark}[Why the dictionary matters]
The hybrid switches between an AESP phase written in terms of \(g(x)\) and a
LocSOR phase often implemented in terms of \(\widehat r(x)\).  The state
handoff must use \cref{eq:gradient-residual-dictionary}; reusing an inner
proximal gradient, or omitting the scalar factor, changes the active set and
invalidates both the certificate and the work bound.
\end{remark}

\subsection{Degree-weighted work}

A local update at vertex \(u\) scans its adjacency list and is charged
\(d_u\) units of work.  For an executed coordinate sequence
\((u_j)_{j=0}^{N-1}\), define
\begin{equation}
    \Work:=\sum_{j=0}^{N-1}d_{u_j}.
    \label{eq:degree-work}
\end{equation}
For a simultaneous active set \(S_k\), one iteration costs
\begin{equation}
    \vol(S_k):=\sum_{u\in S_k}d_u.
    \label{eq:set-volume}
\end{equation}
The nested AESP work at outer stage \(t\) is
\begin{equation}
    \Work_t^{\rm in}
    :=\sum_{k=0}^{K_t-1}\vol(S_t^{(k)}).
    \label{eq:aesp-inner-work}
\end{equation}
Wall-clock time is an important secondary measurement but is not a stable
substitute for \cref{eq:degree-work}: queue implementation, language runtime,
and cache behavior can obscure the scientific locality question.

\subsection{Three quantities that must not be conflated}

The discussion used the symbol \(C_t\) for more than one object.  We henceforth
separate them:
\begin{align}
    F_t
    &:=f(x^{(t)})-f(x^\star),
    &&\text{original objective gap},
    \label{eq:objective-gap-def}\\
    M_t
    &:=\norm{D^{1/2}g(x^{(t)})}_1,
    &&\text{original weighted residual mass},
    \label{eq:mass-def}\\
    \Lambda_t
    &:=\frac{\overline{\vol}(S_t)}{\gamma_t},
    &&\text{AESP locality factor}.
    \label{eq:locality-factor-def}
\end{align}
The desired end-to-end proof needs control of \(\Lambda_t\) during the AESP
burn-in.  A bound on \(M_t\) can instead yield a particularly strong LocSOR
handoff, but the existing AESP theorem does not provide the required
weighted-\(\ell_1\) decay rate.